% Compile beside math-review.sty; replace all visible insertion text for publication.
% Rename and regroup functional headings around the actual mathematical thread.
% Shared files do not determine section hierarchy; no environment/count quota applies.
% Edition check: remove a side branch only when the selected explanation stays intact.
% Copyedit: state a full result once locally; check the printed meaning of notation.
\documentclass[12pt,reqno]{amsart}
\usepackage{math-review}
\title[Part II: Results (concise)]{[Review topic]\\
Part II: Results and Relations (concise)}
\author{[Author name]}
\date{[Prepared date; literature cutoff date]}
\begin{document}
\begin{abstract}
\placeholder{Identify the main question, selected established answers and the mathematical relationships that make them a coherent theory.}
\end{abstract}
\maketitle
\tableofcontents

\section{Introduction}\label{sec:introduction}
\placeholder{State the central question, the perspective used to organize its known answers and the decisive distinctions among those answers. Introduce only the setting needed to make this question readable.}
\placeholder{Explain which part of the question the result group in Section~\ref{sec:results} answers and how the comparisons in Section~\ref{sec:relations} refine that answer. Describe mathematical relationships, not a sequence of paper titles.}

\section{The principal answers to the question}\label{sec:results}
\placeholder{State the necessary objects and conventions, then organize selected results by the question they resolve. Merge technical variants that do not need separate top-level treatment.}
\begin{theorem}[A principal result]\label{thm:principal}
\placeholder{Give the verified complete statement: objects, assumptions, quantifiers, range, conclusion and qualifications. Use the canonical body for a shared full statement, and retain the exact source and controlling version.}
\end{theorem}
\placeholder{Explain the mathematical force of Theorem~\ref{thm:principal}. Give the decisive proof idea: what it changes or constructs, what this achieves and why it leads to the result. A method name alone is insufficient.}

\subsection{The distinction needed to interpret the result}\label{subsec:distinction}
\placeholder{Explain a necessary special case, neighboring conclusion or limitation of scope. Retain conditions while compressing secondary variants; remove this subsection if it serves no actual distinction.}

\section{Relationships and consequences}\label{sec:relations}
\placeholder{Compare the chosen answers by matching assumptions and conclusions. Prove or cite a claimed implication, improvement or equivalence; distinguish it from analogy. Derive the principal consequence and return to what the central question has and has not been answered within the selected scope.}

\templatereferences
\end{document}
